\documentclass[12pt, a4paper]{ctexart}

%==================== 导言区：宏包加载 ====================
\usepackage{amsmath, amssymb, amsthm} % 数学公式与符号
\usepackage{geometry}                 % 页面设置
\usepackage{xcolor}                   % 颜色支持
\usepackage{fancyhdr}                 % 页眉页脚
\usepackage{enumitem}                 % 列表定制
\usepackage{tcolorbox}                %以此美化文本框
\usepackage{titlesec}                 % 标题格式调整

%==================== 页面布局设置 ====================
\geometry{left=2.5cm, right=2.5cm, top=2.5cm, bottom=2.5cm}

%==================== 自定义视觉样式 ====================

% 1. 定义分隔线样式
\newcommand{\TopSeparator}{
    \par\noindent\rule{\textwidth}{1.5pt}\vspace{0.5em} % 粗实线 ━━━━
}
\newcommand{\AnalysisSeparator}{
    \par\noindent\rule{\textwidth}{0.5pt}\vspace{0.2em}\hrule\vspace{0.5em} % 双线效果 ════
}

% 2. 定义红色解析环境
%在这个环境内的所有文字和公式都会自动变成红色
\newenvironment{RedAnalysis}{
    \par
    \color{red} % 设置后续字体为红色
}{
    \par
    \color{black} % 结束时恢复黑色
}

% 3. 标题格式微调
\titleformat{\section}{\large\bfseries\heiti}{}{0em}{}[\vspace{-0.5em}]
\titleformat{\subsection}{\bfseries\kaishu}{}{0em}{}

\newcommand{\choicesFour}[4]{
    \begin{enumerate}[label=\Alph*., itemsep=0pt, parsep=0pt, topsep=5pt, labelsep=0.5em]
        \begin{minipage}{0.22\linewidth} \item #1 \end{minipage}
        \begin{minipage}{0.22\linewidth} \item #2 \end{minipage}
        \begin{minipage}{0.22\linewidth} \item #3 \end{minipage}
        \begin{minipage}{0.22\linewidth} \item #4 \end{minipage}
    \end{enumerate}
}
\newcommand{\choicesTwo}[4]{
    \begin{enumerate}[label=\Alph*., itemsep=0pt, parsep=0pt, topsep=5pt, labelsep=0.5em]
        \begin{minipage}{0.45\linewidth} \item #1 \end{minipage}
        \begin{minipage}{0.45\linewidth} \item #2 \end{minipage}
        \begin{minipage}{0.45\linewidth} \item #3 \end{minipage}
        \begin{minipage}{0.45\linewidth} \item #4 \end{minipage}
    \end{enumerate}
}
\newcommand{\choices}[4]{
    \begin{enumerate}[label=\Alph*., itemsep=0pt, parsep=0pt, topsep=5pt, labelsep=0.5em]
        \item #1
        \item #2
        \item #3
        \item #4
    \end{enumerate}
}

%==================== 文档开始 ====================
\begin{document}

% 标题信息
\title{\textbf{第十一章 二重积分（解析）}}
\author{数学习题讲义}
\date{\today}
\maketitle

% --- 题目 1 ---
\TopSeparator
\section*{题目1}

\textbf{【题目重现】} \\
设区域 $ D = \{(x, y) \mid 1 \leq x^{2} + y^{2} \leq 2\} $ ，则二重积分 $ \iint_{D} dx dy = $ \underline{\hspace{1cm}}。
\choicesFour{$ \pi $}{$ 2\pi $}{$ 3\pi $}{$ 4\pi $}

\AnalysisSeparator

\begin{RedAnalysis}

\subsection*{一、逐步解析}

\textbf{第一步：问题分析} \\
被积函数为1，二重积分的几何意义是求区域 $D$ 的面积。
$D$ 是圆环，内半径 $R_{in} = 1$，外半径 $R_{out} = \sqrt{2}$。

\textbf{详细步骤} \\
1. 面积 $ \sigma(D) = \pi (\sqrt{2})^2 - \pi (1)^2 $。
2. 计算：$ \pi (2) - \pi (1) = \pi $。

\textbf{第四步：最终答案} \\
A

\subsection*{二、核心公式}
1. \textbf{几何意义}：$\iint_D 1 dxdy = S_D$。

\end{RedAnalysis}

\vspace{2em}

% --- 题目 2 ---
\TopSeparator
\section*{题目2}

\textbf{【题目重现】} \\
已知 $ F(x, y) = \ln(1 + x^{2} + y^{2}) + \iint_{D} \, dx \, dy $ , 其中 $D$ 为 $xOy$ 坐标平面上的有界闭区域且 $ f(x, y) $ 在 $D$ 上连续，则 $ F(x, y) $ 在点 $(1, 2)$ 处的全微分为 \underline{\hspace{1cm}}。
\choicesTwo
{$ \frac{1}{3}dx+\frac{2}{3}dy $}
{$ \frac{1}{3}dx+\frac{2}{3}dy+f(1,2) $}
{$ \frac{2}{3}dx+\frac{1}{3}dy $}
{$ \frac{2}{3}dx+\frac{1}{3}dy+f(1,2) $}

\AnalysisSeparator

\begin{RedAnalysis}

\subsection*{一、逐步解析}

\textbf{第一步：问题分析} \\
注意 $\iint_D dxdy$ 是常数（即区域 $D$ 的面积），对常数微分结果为0。
所以只需对第一项 $\ln(1+x^2+y^2)$ 求微分。

\textbf{详细步骤} \\
1. 令 $u = \ln(1+x^2+y^2)$。
2. $u_x = \frac{2x}{1+x^2+y^2}, u_y = \frac{2y}{1+x^2+y^2}$。
3. 在 $(1,2)$ 处：
   $1+x^2+y^2 = 1+1+4 = 6$。
   $u_x = 2/6 = 1/3$。
   $u_y = 4/6 = 2/3$。
4. 全微分 $dF = (1/3)dx + (2/3)dy$。

\textbf{第四步：最终答案} \\
A

\subsection*{二、核心公式}
1. \textbf{常数导数为0}。

\end{RedAnalysis}

\vspace{2em}

% --- 题目 3 ---
\TopSeparator
\section*{题目3}

\textbf{【题目重现】} \\
如果闭区域 $D$ 由 $x$ 轴，$y$ 轴及 $ x+y=1 $ 围成。则 $ \iint_{D}(x+y)^{2}d\sigma $ \underline{\hspace{1cm}} $ \iint_{D}(x+y)^{3}d\sigma $。
\choicesFour{大于}{小于}{等于}{不确定}

\AnalysisSeparator

\begin{RedAnalysis}

\subsection*{一、逐步解析}

\textbf{第一步：问题分析} \\
比较积分大小。积分性质：若在 $D$ 上 $f(x,y) > g(x,y)$，则 $\iint_D f > \iint_D g$。
分析定义域上的值域。

\textbf{详细步骤} \\
1. 区域 $D$ 满足 $0 \le x+y \le 1$（三角形区域）。
2. 令 $t = x+y$，则 $0 \le t \le 1$。
3. 当 $0 < t < 1$ 时，$t^2 > t^3$（因为 $t$ 越乘越小）。
4. 所以在积分区域（除边界）上，$(x+y)^2 > (x+y)^3$。
5. 积分具有保号性，故前者积分大于后者。

\textbf{第四步：最终答案} \\
A

\subsection*{二、核心公式}
1. \textbf{比较性质}：$f \ge g \Rightarrow \iint f \ge \iint g$。

\end{RedAnalysis}

\vspace{2em}

% --- 题目 4 ---
\TopSeparator
\section*{题目4}

\textbf{【题目重现】} \\
设平面区域 $ D=\{(x,y)|x^{2}+y^{2}\leq1\} $ ， $ D_{1}=\{(x,y)|x^{2}+y^{2}\leq1,x\geq0,y\geq0\} $ 则下列等式一定成立的是 \underline{\hspace{1cm}}。
\choicesTwo
{$ \iint_{D}\sqrt{x^{2}+y^{2}}dxdy=4\iint_{D_{1}}\sqrt{x^{2}+y^{2}}dxdy $}
{$ \iint_{D} xydxdy = 4 \iint_{D_{1}} xydxdy $}
{$ \iint_{D} f(x,y)dxdy = 4 \iint_{D_{1}} f(x,y)dxdy $}
{$ \iint_{D} xdxdy = 4 \iint_{D_{1}} xdxdy $}

\AnalysisSeparator

\begin{RedAnalysis}

\subsection*{一、逐步解析}

\textbf{第一步：问题分析} \\
利用对称性。$D$ 是单位圆，$D_1$ 是第一象限部分（1/4圆）。
若 $f(x,y)$ 关于 $x, y$ 均为偶函数（即关于坐标轴对称），则积分为 4 倍。

\textbf{详细步骤} \\
1. **选项A**：$f = \sqrt{x^2+y^2}$。
   $f(-x,y) = f(x,y)$，$f(x,-y) = f(x,y)$。
   关于 $x$ 轴、$y$ 轴对称。
   所以总积分 = 4倍第一象限积分。正确。
2. **选项B**：$xy$。奇函数（关于 $x$ 和关于 $y$）。积分为0。不等于4倍。
3. **选项C**：一般函数不一定对称。
4. **选项D**：$x$。奇函数（关于 $x$ 轴不对称，关于 $y$ 轴奇函数），积分为0。

\textbf{第四步：最终答案} \\
A

\subsection*{二、核心公式}
1. \textbf{对称性}：偶函数在对称区域积分为倍数，奇函数为0。

\end{RedAnalysis}

\vspace{2em}

% --- 题目 5 ---
\TopSeparator
\section*{题目5}

\textbf{【题目重现】} \\
积分区域 $D$ 为 $ x^{2}+y^{2}\leq2 $ , 则 $ \iint_{D}xd\sigma= $ \underline{\hspace{1cm}}。
\choicesFour{$ 2\pi $}{$ \pi $}{1}{0}

\AnalysisSeparator

\begin{RedAnalysis}

\subsection*{一、逐步解析}

\textbf{第一步：问题分析} \\
利用奇偶对称性。

\textbf{详细步骤} \\
1. 区域 $D$ 关于 $y$ 轴对称（圆）。
2. 被积函数 $f(x,y) = x$ 是关于 $x$ 的奇函数（$f(-x,y) = -x = -f(x,y)$）。
3. 根据奇函数在对称区域积分为0的性质，结果为0。

\textbf{第四步：最终答案} \\
D

\subsection*{二、核心公式}
1. \textbf{对称性}：$\iint_D x dxdy = 0$ ($D$关于$y$轴对称)。

\end{RedAnalysis}

\vspace{2em}

% --- 题目 6 ---
\TopSeparator
\section*{题目6}

\textbf{【题目重现】} \\
$D$ 是以原点为圆心，半径为 5 的圆形区域。则 $ \iint_{D} d\sigma $ \underline{\hspace{2cm}}。

\AnalysisSeparator

\begin{RedAnalysis}

\subsection*{一、逐步解析}

\textbf{详细步骤} \\
1. 几何意义是求面积。
2. $S = \pi r^2 = \pi (5)^2 = 25\pi$。

\textbf{第四步：最终答案} \\
$ 25\pi $

\end{RedAnalysis}

\vspace{2em}

% --- 题目 7 ---
\TopSeparator
\section*{题目7}

\textbf{【题目重现】} \\
设区域 $D$ 是由直线 $y = 1, x = 2$ 及 $y = x$ 所围成的三角形闭区域，则二重积分 $ \iint_{D} xy \, dx \, dy = $ \underline{\hspace{2cm}}。

\AnalysisSeparator

\begin{RedAnalysis}

\subsection*{一、逐步解析}

\textbf{第一步：问题分析} \\
确定积分限，化为累次积分。
区域：$y=1$（下），$y=x$（上/左），$x=2$（右）。
交点：$y=1$ 与 $y=x \Rightarrow (1,1)$；$x=2$ 与 $y=1 \Rightarrow (2,1)$；$x=2$ 与 $y=x \Rightarrow (2,2)$。
所以 $1 \le x \le 2$，且 $1 \le y \le x$。

\textbf{详细步骤} \\
1. 写成X型区域：$1 \le x \le 2, 1 \le y \le x$。
2. 积分：
   \[ \int_1^2 dx \int_1^x xy dy = \int_1^2 x [\frac{1}{2}y^2]_1^x dx \]
   \[ = \int_1^2 \frac{1}{2}x (x^2 - 1) dx = \frac{1}{2} \int_1^2 (x^3 - x) dx \]
   \[ = \frac{1}{2} [\frac{1}{4}x^4 - \frac{1}{2}x^2]_1^2 \]
   \[ = \frac{1}{2} [ (4 - 2) - (\frac{1}{4} - \frac{1}{2}) ] = \frac{1}{2} [ 2 - (-\frac{1}{4}) ] = \frac{1}{2} \cdot \frac{9}{4} = \frac{9}{8} \]

\textbf{第四步：最终答案} \\
$ \frac{9}{8} $

\subsection*{二、核心公式}
1. \textbf{累次积分}：$\iint_D f(x,y) dxdy = \int_a^b dx \int_{\phi_1(x)}^{\phi_2(x)} f(x,y) dy$。

\end{RedAnalysis}

\vspace{2em}

% --- 题目 8 ---
\TopSeparator
\section*{题目8}

\textbf{【题目重现】} \\
区域 $D$ 是 $ x^{2} + y^{2} \leq a^{2} $ ，计算 $ \iint_{D} e^{-x^{2}-y^{2}} \, dx \, dy = $ \underline{\hspace{2cm}}。

\AnalysisSeparator

\begin{RedAnalysis}

\subsection*{一、逐步解析}

\textbf{第一步：问题分析} \\
圆域，被积函数含 $x^2+y^2$，使用极坐标。

\textbf{详细步骤} \\
1. 极坐标变换：$x^2+y^2 = r^2, dxdy = r dr d\theta$。
2. 范围：$0 \le \theta \le 2\pi, 0 \le r \le a$。
3. 积分：
   \[ \int_0^{2\pi} d\theta \int_0^a e^{-r^2} r dr \]
   \[ = 2\pi \cdot [-\frac{1}{2}e^{-r^2}]_0^a \]
   \[ = 2\pi (-\frac{1}{2})(e^{-a^2} - 1) = \pi (1 - e^{-a^2}) \]

\textbf{第四步：最终答案} \\
$ \pi (1 - e^{-a^2}) $

\subsection*{二、核心公式}
1. \textbf{极坐标积分}：$dxdy = r dr d\theta$。

\end{RedAnalysis}

\vspace{2em}

% --- 题目 9 ---
\TopSeparator
\section*{题目9}

\textbf{【题目重现】} \\
$ D = \{(x, y) \mid 1 \leq x^{2} + y^{2} \leq 4\} $ 求 $ \iint_{D} \ln(x^{2} + y^{2}) \, dx \, dy = $ \underline{\hspace{2cm}}。

\AnalysisSeparator

\begin{RedAnalysis}

\subsection*{一、逐步解析}

\textbf{第一步：问题分析} \\
圆环域，极坐标。$1 \le r \le 2$。

\textbf{详细步骤} \\
1. 令 $x^2+y^2 = r^2$。
2. 积分：
   \[ \int_0^{2\pi} d\theta \int_1^2 \ln(r^2) \cdot r dr \]
   \[ = 2\pi \int_1^2 2r \ln r dr = 4\pi \int_1^2 r \ln r dr \]
3. 分部积分 $\int r \ln r dr$：
   令 $u=\ln r, dv=r dr \Rightarrow v=r^2/2$。
   $[\frac{r^2}{2}\ln r]_1^2 - \int_1^2 \frac{r^2}{2} \cdot \frac{1}{r} dr$
   $= (2\ln 2 - 0) - \int_1^2 \frac{r}{2} dr$
   $= 2\ln 2 - [\frac{r^2}{4}]_1^2 = 2\ln 2 - (1 - 1/4) = 2\ln 2 - 3/4$。
4. 总结果：
   $4\pi (2\ln 2 - 3/4) = 8\pi \ln 2 - 3\pi$。

\textbf{第四步：最终答案} \\
$ 8\pi \ln 2 - 3\pi $

\end{RedAnalysis}

\vspace{2em}

% --- 题目 10 ---
\TopSeparator
\section*{题目10}

\textbf{【题目重现】} \\
设 $D$ 是矩形： $ 0 \leq x \leq a, 0 \leq y \leq b $ ，则 $ \iint_{D} dxdy = $ \underline{\hspace{2cm}}。

\AnalysisSeparator

\begin{RedAnalysis}

\subsection*{一、逐步解析}

\textbf{详细步骤} \\
求矩形面积。
$S = a \cdot b = ab$。

\textbf{第四步：最终答案} \\
$ ab $

\end{RedAnalysis}

\vspace{2em}

% --- 题目 11 ---
\TopSeparator
\section*{题目11}

\textbf{【题目重现】} \\
求二重积分 $ \iint_{D}(2x+y)dxdy $ , 其中 $D$ 是由直线 $y=3, y=x$ 及曲线 $xy=1$ 围成的闭区域。

\AnalysisSeparator

\begin{RedAnalysis}

\subsection*{一、逐步解析}

\textbf{第一步：问题分析} \\
画图确定区域与积分次序。
$y=x$ 与 $xy=1 \Rightarrow x^2=1 \Rightarrow (1,1)$。
$y=3$ 与 $y=x \Rightarrow (3,3)$。
$y=3$ 与 $xy=1 \Rightarrow (1/3, 3)$。
区域在 $y=1$ 到 $y=3$ 之间比较好积分（Y型）。
左边界：$x = 1/y$（双曲线）。
右边界：$x = y$（直线）。
范围：$1 \le y \le 3, 1/y \le x \le y$。

\textbf{详细步骤} \\
1. 设置积分：
   \[ \int_1^3 dy \int_{1/y}^y (2x+y) dx \]
2. 对 $x$ 积分：
   \[ [x^2 + xy]_{1/y}^y = (y^2 + y^2) - (\frac{1}{y^2} + 1) = 2y^2 - y^{-2} - 1 \]
3. 对 $y$ 积分：
   \[ \int_1^3 (2y^2 - y^{-2} - 1) dy = [\frac{2}{3}y^3 + y^{-1} - y]_1^3 \]
   \[ = (\frac{54}{3} + \frac{1}{3} - 3) - (\frac{2}{3} + 1 - 1) \]
   \[ = (18 + \frac{1}{3} - 3) - \frac{2}{3} = 15 - \frac{1}{3} = \frac{44}{3} \]

\textbf{第四步：最终答案} \\
$ \frac{44}{3} $

\subsection*{二、核心公式}
1. \textbf{Y型区域积分}。

\end{RedAnalysis}

\vspace{2em}

% --- 题目 12 ---
\TopSeparator
\section*{题目12}

\textbf{【题目重现】} \\
计算二重积分 $ \iint_{D} xy \, dx \, dy $ , 其中 $D$ 是由直线 $ y = x, y = 5x $ 及 $ y = -x + 6 $ 所围成的闭区域。

\AnalysisSeparator

\begin{RedAnalysis}

\subsection*{一、逐步解析}

\textbf{第一步：问题分析} \\
找交点：
$y=x$ 与 $y=-x+6 \Rightarrow 2x=6 \Rightarrow (3,3)$。
$y=5x$ 与 $y=-x+6 \Rightarrow 6x=6 \Rightarrow (1,5)$。
$y=x$ 与 $y=5x \Rightarrow (0,0)$。
三角形区域。最好分块或选合适次序。
用X型：分段点 $x=1$。
$0 \le x \le 1$: $x \le y \le 5x$。
$1 \le x \le 3$: $x \le y \le 6-x$。

\textbf{详细步骤} \\
1. 第一路：
   \[ \int_0^1 x dx \int_x^{5x} y dy + \int_1^3 x dx \int_x^{6-x} y dy \]
2. 计算第一部分：
   $\int_x^{5x} y dy = \frac{1}{2}(25x^2 - x^2) = 12x^2$。
   $\int_0^1 x(12x^2) dx = 3[x^4]_0^1 = 3$。
3. 计算第二部分：
   $\int_x^{6-x} y dy = \frac{1}{2}((6-x)^2 - x^2) = \frac{1}{2}(36 - 12x) = 18 - 6x$。
   $\int_1^3 x(18 - 6x) dx = \int_1^3 (18x - 6x^2) dx = [9x^2 - 2x^3]_1^3$。
   $= (81 - 54) - (9 - 2) = 27 - 7 = 20$。
4. 总和：$3 + 20 = 23$。

\textbf{第四步：最终答案} \\
23

\end{RedAnalysis}

\vspace{2em}

% --- 题目 13 ---
\TopSeparator
\section*{题目13}

\textbf{【题目重现】} \\
计算二重积分 $ \iint_{D} xydxdy $ , 其中 $D$ 为直线 $ y = 1, x = 2 $ 及 $ y = x $ 所围成的闭区域。

\AnalysisSeparator

\begin{RedAnalysis}

\subsection*{一、逐步解析}

\textbf{详细步骤} \\
同第7题，题目完全一致。
答案：$9/8$。

\textbf{第四步：最终答案} \\
$ \frac{9}{8} $

\end{RedAnalysis}

\vspace{2em}

% --- 题目 14 ---
\TopSeparator
\section*{题目14}

\textbf{【题目重现】} \\
求二重积分 $ \iint_{D} x^{2} dx dy $ , 其中 $D$ 是由 $ x = 0, y = 2x, y = 2 $ 所围成的区域。

\AnalysisSeparator

\begin{RedAnalysis}

\subsection*{一、逐步解析}

\textbf{第一步：问题分析} \\
区域形状：左 $x=0$，上 $y=2$，斜 $y=2x$。
交点 $(0,0), (0,2), (1,2)$。
选择Y型积分较简单：$0 \le y \le 2, 0 \le x \le y/2$。

\textbf{详细步骤} \\
1. 设置积分：
   \[ \int_0^2 dy \int_0^{y/2} x^2 dx \]
2. 对 $x$ 积分：
   \[ \frac{1}{3} [(y/2)^3 - 0] = \frac{1}{24}y^3 \]
3. 对 $y$ 积分：
   \[ \int_0^2 \frac{1}{24}y^3 dy = \frac{1}{24} [\frac{1}{4}y^4]_0^2 \]
   \[ = \frac{1}{96} (16) = \frac{1}{6} \]

\textbf{第四步：最终答案} \\
$ \frac{1}{6} $

\end{RedAnalysis}

\vspace{2em}

% --- 题目 15 ---
\TopSeparator
\section*{题目15}

\textbf{【题目重现】} \\
求 $ I = \iint_{D} (3x + 2y) d\sigma $ ，其中 $D$ 是由两坐标轴及直线 $ x + y = 2 $ 所围成的闭区域。

\AnalysisSeparator

\begin{RedAnalysis}

\subsection*{一、逐步解析}

\textbf{第一步：问题分析} \\
三角形区域：$0 \le x \le 2, 0 \le y \le 2-x$。

\textbf{详细步骤} \\
1. 设置积分：
   \[ \int_0^2 dx \int_0^{2-x} (3x + 2y) dy \]
2. 对 $y$ 积分：
   \[ [3xy + y^2]_0^{2-x} = 3x(2-x) + (2-x)^2 = 6x - 3x^2 + 4 - 4x + x^2 = -2x^2 + 2x + 4 \]
3. 对 $x$ 积分：
   \[ \int_0^2 (-2x^2 + 2x + 4) dx = [-\frac{2}{3}x^3 + x^2 + 4x]_0^2 \]
   \[ = -\frac{16}{3} + 4 + 8 = 12 - \frac{16}{3} = \frac{36-16}{3} = \frac{20}{3} \]

\textbf{第四步：最终答案} \\
$ \frac{20}{3} $

\end{RedAnalysis}

\end{document}
